% answers_short.tex —— 2012 年·全国大纲卷·理科 · 简短答案速查
% 由 tools/make_short_answers.py 自 answers.tex 的 \gksolhead/\gkanswer 提取，
% 并与 2012/2012_数学真题_全国大纲理_答案要点.md 逐题比对后生成，勿手改（改源 MD 后重跑）。

\begin{gknotice}
\textbf{使用说明}：本节只给最终答案便于快速核对，完整推导见解析版。答案为编写者独立推导，\textbf{非官方标准答案}。
\end{gknotice}

\gkpart{选择题}
\begin{center}\begin{tabular}{ccccccccccccc}
\toprule
题号 & 1 & 2 & 3 & 4 & 5 & 6 & 7 & 8 & 9 & 10 & 11 & 12 \\
\midrule
答案 & C & B & C & D & A & D & A & A & D & A & A & B \\
\bottomrule
\end{tabular}\end{center}

\gkpart{填空题}
\begin{itemize}[leftmargin=2.6em]
  \item[\textbf{13.}] $-1$
  \item[\textbf{14.}] $\dfrac{5\pi}{6}$
  \item[\textbf{15.}] $56$
  \item[\textbf{16.}] $\dfrac{\sqrt6}{6}$
\end{itemize}

\gkpart{解答题}
\begin{itemize}[leftmargin=2.6em]
  \item[\textbf{17.}] $C=\dfrac\pi6$。
  \item[\textbf{18.}] (1) 见解析版；(2) $PD$ 与平面 $PBC$ 所成角为 $30^\circ$
  \item[\textbf{19.}] (1) $\dfrac{44}{125}$；(2) $E(\xi)=\dfrac75$
  \item[\textbf{20.}] (1) 当 $a\leqslant0$ 时，$f$ 在 $[0,\pi]$ 上单调递减；当 $0<a<1$ 时，令 $\alpha=\arcsin a$，则 $f$ 在 $[0,\alpha]$、$[\pi-\alpha,\pi]$ 上递增，在 $[\alpha,\pi-\alpha]$ 上递减；当 $a\geqslant1$ 时，$f$ 在 $[0,\pi]$ 上单调递增；(2) $a\leqslant\dfrac2\pi$
  \item[\textbf{21.}] (1) $r=\dfrac{\sqrt5}{2}$；(2) $\dfrac6{\sqrt5}=\dfrac{6\sqrt5}{5}$
  \item[\textbf{22.}] (1) 见解析版；(2) $x_n=\dfrac{9\cdot5^{n-1}-1}{3\cdot5^{n-1}+1}$
\end{itemize}
